\documentclass[11pt]{article}
\usepackage{amsmath,amssymb,amsthm}
\usepackage[margin=1in]{geometry}

\newtheorem{theorem}{Theorem}
\newtheorem{lemma}[theorem]{Lemma}
\newtheorem{proposition}[theorem]{Proposition}
\newtheorem{corollary}[theorem]{Corollary}

\title{Solution to Ramanujan Challenge Problem 2.1:\\
Polynomial Continued Fraction for $6/(3-\pi)$}
\author{Xiang Huang}
\date{July 2026}

\begin{document}
\maketitle

\begin{abstract}
We prove that the polynomial continued fraction (PCF) with
$a_n = -220n^3 - 484n^2 - 301n - 42$ and
$b_n = 4n^2(2n+1)^2(5n-4)(5n+6)$ converges to $6/(3-\pi)$.
The proof reduces to a known continued fraction for~$\pi$ from
Cohen's database (arXiv:2409.06086, Entry~1.2.18) via an
elementary sign-flip identity on convergents.
\end{abstract}

\section{Statement}

Define
\begin{align}
a_n &= -220n^3 - 484n^2 - 301n - 42, \label{eq:an}\\
b_n &= 4n^2(2n+1)^2(5n-4)(5n+6). \label{eq:bn}
\end{align}
We prove:
\begin{equation}\label{eq:main}
a_0 + \cfrac{b_1}{a_1 + \cfrac{b_2}{a_2 + \cfrac{b_3}{a_3 + \cdots}}}
= \frac{6}{3-\pi}.
\end{equation}

\section{Cohen's continued fraction for $\pi$}

Henri Cohen's database of polynomial continued fractions
\cite{Cohen2024} contains Entry~1.2.18:
\begin{equation}\label{eq:cohen}
\pi = 3 + \cfrac{6}{\alpha(1) + \cfrac{\beta(1)}{\alpha(2) +
\cfrac{\beta(2)}{\alpha(3) + \cdots}}}
\end{equation}
where, for $n \ge 1$,
\begin{align}
\alpha(n) &= 220n^3 - 176n^2 - 7n + 5, \label{eq:alpha}\\
\beta(n) &= 4n^2(2n+1)^2(5n-4)(5n+6). \label{eq:beta}
\end{align}
This continued fraction converges to~$\pi$ at an exponential
rate of $\varphi^{-10} \approx 10^{-2.09}$ digits per term,
where $\varphi = (1+\sqrt{5})/2$ is the golden ratio.
We have verified~\eqref{eq:cohen} numerically to over 1000
decimal digits.

\section{Polynomial identities}

\begin{lemma}[Index shift]\label{lem:shift}
For all $n \ge 0$,
\[
a_n = -\alpha(n+1) \quad\text{and}\quad b_n = \beta(n).
\]
\end{lemma}

\begin{proof}
For the first identity:
\begin{align*}
\alpha(n+1) &= 220(n+1)^3 - 176(n+1)^2 - 7(n+1) + 5 \\
&= 220n^3 + 660n^2 + 660n + 220 - 176n^2 - 352n - 176 - 7n - 7 + 5 \\
&= 220n^3 + 484n^2 + 301n + 42.
\end{align*}
Hence $-\alpha(n+1) = -220n^3 - 484n^2 - 301n - 42 = a_n$.

The second identity is immediate: both $b_n$ and $\beta(n)$ are
$4n^2(2n+1)^2(5n-4)(5n+6)$. \qedhere
\end{proof}

\section{Sign-flip lemma}

\begin{lemma}[Sign-flip identity for convergents]\label{lem:signflip}
Let $\{c_k\}_{k \ge 0}$ and $\{d_k\}_{k \ge 1}$ be sequences
with $d_k \ne 0$, and define the convergents of the continued
fraction $S = c_0 + d_1/(c_1 + d_2/(c_2 + \cdots))$ by the
recurrence
\begin{equation}\label{eq:recurrence}
X_{n} = c_n X_{n-1} + d_n X_{n-2}, \quad n \ge 1,
\end{equation}
with $P_{-1} = 1$, $P_0 = c_0$, $Q_{-1} = 0$, $Q_0 = 1$.

Define the ``sign-flipped'' continued fraction
$\widetilde{S} = -c_0 + d_1/(-c_1 + d_2/(-c_2 + \cdots))$,
whose convergents satisfy
\begin{equation}\label{eq:recurrence2}
\widetilde{X}_{n} = -c_n \widetilde{X}_{n-1} + d_n \widetilde{X}_{n-2},
\quad n \ge 1,
\end{equation}
with $\widetilde{P}_{-1} = 1$, $\widetilde{P}_0 = -c_0$,
$\widetilde{Q}_{-1} = 0$, $\widetilde{Q}_0 = 1$.

Then for all $n \ge -1$:
\begin{equation}\label{eq:signrel}
\widetilde{P}_n = (-1)^{n+1} P_n
\quad\text{and}\quad
\widetilde{Q}_n = (-1)^n Q_n.
\end{equation}
In particular, $\widetilde{P}_n/\widetilde{Q}_n = -P_n/Q_n$
for all $n \ge 0$, so
\begin{equation}\label{eq:negconv}
\widetilde{S} = -S
\end{equation}
whenever $S$ converges.
\end{lemma}

\begin{proof}
We prove~\eqref{eq:signrel} by induction on $n$.

\emph{Base cases.}
$\widetilde{P}_{-1} = 1 = (-1)^0 \cdot 1 = (-1)^0 P_{-1}$. \checkmark

$\widetilde{P}_0 = -c_0 = (-1)^1 c_0 = (-1)^1 P_0$. \checkmark

$\widetilde{Q}_{-1} = 0 = (-1)^{-1} \cdot 0 = (-1)^{-1} Q_{-1}$. \checkmark

$\widetilde{Q}_0 = 1 = (-1)^0 \cdot 1 = (-1)^0 Q_0$. \checkmark

\emph{Inductive step.} Suppose $\widetilde{P}_{n-1} = (-1)^n P_{n-1}$
and $\widetilde{P}_{n-2} = (-1)^{n-1} P_{n-2}$. Then:
\begin{align*}
\widetilde{P}_n &= (-c_n)\widetilde{P}_{n-1} + d_n\widetilde{P}_{n-2} \\
&= (-c_n)(-1)^n P_{n-1} + d_n (-1)^{n-1} P_{n-2} \\
&= (-1)^{n+1} c_n P_{n-1} + (-1)^{n+1}(-1)^{-2} d_n P_{n-2} \\
&= (-1)^{n+1} (c_n P_{n-1} + d_n P_{n-2}) \\
&= (-1)^{n+1} P_n.
\end{align*}
In the third line we used $(-1)^{n-1} = (-1)^{n+1}$ since
$(-1)^{-2} = 1$. The induction for $\widetilde{Q}_n$ is
identical (with $\widetilde{Q}_{-1} = Q_{-1} = 0$). \qedhere
\end{proof}

\section{Proof of the main identity}

\begin{theorem}\label{thm:main}
The continued fraction~\eqref{eq:main} converges to $6/(3-\pi)$.
\end{theorem}

\begin{proof}
Let $T$ denote the first tail of Cohen's continued
fraction~\eqref{eq:cohen}:
\[
T = \alpha(1) + \cfrac{\beta(1)}{\alpha(2) +
\cfrac{\beta(2)}{\alpha(3) + \cdots}}.
\]
From~\eqref{eq:cohen}, $\pi = 3 + 6/T$, hence
\begin{equation}\label{eq:T}
T = \frac{6}{\pi - 3}.
\end{equation}

By Lemma~\ref{lem:shift}, our challenge PCF has partial
quotients $a_k = -\alpha(k+1)$ and partial numerators
$b_k = \beta(k)$ for $k \ge 1$, so it equals
\[
-\alpha(1) + \cfrac{\beta(1)}{-\alpha(2) +
\cfrac{\beta(2)}{-\alpha(3) + \cdots}}.
\]
This is exactly the sign-flipped version of $T$
(Definition of Lemma~\ref{lem:signflip}, with $c_k = \alpha(k+1)$
and $d_k = \beta(k)$). By~\eqref{eq:negconv},
\[
a_0 + \cfrac{b_1}{a_1 + \cfrac{b_2}{a_2 + \cdots}}
= -T = -\frac{6}{\pi - 3} = \frac{6}{3 - \pi}.
\qedhere
\]
\end{proof}

\section{Numerical verification}

We have independently verified the identity~\eqref{eq:main}
and Cohen's base identity~\eqref{eq:cohen} to over 1000
decimal digits using arbitrary-precision arithmetic (Python
\texttt{mpmath}). With $N$ terms computed:
\begin{itemize}
\item Cohen's CF~\eqref{eq:cohen} matches $\pi$ to
  $\lfloor 2.09 N \rfloor$ digits, consistent with Poincar\'e
  roots $\varphi^5$ and $-\varphi^{-5}$ (convergence base
  $\varphi^{-10}$).
\item Our PCF~\eqref{eq:main} converges at the same rate,
  as the sign-flip lemma preserves convergent magnitude.
\end{itemize}

\section{Remarks}

\begin{enumerate}
\item The Poincar\'e polynomial of the three-term
  recurrence~\eqref{eq:recurrence} is $c^2 - 11c - 1 = 0$,
  with roots $\varphi^5 = (11 + 5\sqrt{5})/2$ and
  $-\varphi^{-5} = (-11 + 5\sqrt{5})/2$.

\item Cohen's CF is listed as Entry~1.2.18 in the 2024
  database~\cite{Cohen2024} and as Entry~5.3.22 in the 2026
  dictionary~\cite{Cohen2026}. The 2026 dictionary identifies
  it as the even contraction of the period-two
  fraction~5.3.16.

\item The partial numerator polynomial $\beta(n) =
  4n^2(2n+1)^2(5n-4)(5n+6)$ factors into
  $\binom{2n}{n}$-related Pochhammer products. The cubic
  Pochhammer gauge $g_n = \frac{(6/5)_n (4/5)_n (1)_n}{(n!)^2}$
  transforms the continuant recurrence into a degree-3
  Ap\'ery-type recurrence (cf.~\cite{Cohen2026}, Section~5.3).
\end{enumerate}

\begin{thebibliography}{10}
\bibitem{Cohen2024}
H.~Cohen,
\emph{A Database of Continued Fractions of Polynomial Type},
arXiv:2409.06086, 2024.

\bibitem{Cohen2026}
H.~Cohen,
\emph{Continued Fractions of Polynomial Type: Theory and
  Encyclopedic Dictionary},
arXiv:2607.06581, 2026.

\bibitem{Challenge2026}
Ramanujan Machine Group,
\emph{The Ramanujan Challenge for AI}, 2026.
\end{thebibliography}

\end{document}
